\chapter{La fachada canónica de API}
\label{ch:api-facade}

\begin{note}
Stub. \code{GraphHunterApi} es el único objeto del que cada
transport (Tauri, HTTP, MCP) mantiene un \code{Arc}. Toda la
gestión de sesión, cache de hunt, conversación AI, streaming
Sentinel, scorer GNN y estado GeoIP vive dentro de \code{ApiState}.
\end{note}

\section{Estado}

\begin{lstlisting}[style=rust]
pub struct GraphHunterApi(Arc<ApiState>);

struct ApiState {
    sessions:           SessionStore,
    hunt_cache:         Arc<RwLock<Vec<Vec<String>>>>,
    npu_scorer:         Arc<RwLock<Option<NpuScorer>>>,
    geoip_provider:     Arc<Mutex<Option<MaxMindProvider>>>,
    http_geoip_client:  Arc<Mutex<Option<Arc<dyn HttpClient>>>>,
    sentinel_connector: Arc<RwLock<Option<SentinelConnectorHandle>>>,
    ai_config:          Arc<RwLock<Option<ProviderConfig>>>,
    emitter:            Arc<dyn EventEmitter>,
    // ... 8 campos mas
}
\end{lstlisting}

La fachada la introdujo el refactor P1. Antes, cada transport
sostenía su propio estado con forma distinta; con la fachada en
su lugar, Tauri / HTTP / MCP son estrictamente capas de
delegación.

\section{Módulo de operaciones}

Cada método de \code{GraphHunterApi} se define en
\filepath{graph\_hunter\_api/src/operations/*.rs}, un archivo por
categoría (analytics, anomaly, dataset, dsl, enrichment,
entity, export, geoip, graph\_ops, hunt, ingestion, sentinel,
session, sigma, ticket, ai). Un método acepta un request DTO,
despacha la sesión si hace falta, llama al motor y devuelve un
response DTO.

\section{Handle de sesión}

\code{SessionHandle(String)} es un ID opaco de sesión; los
requests o lo llevan (target explícito) o lo omiten (usar la
sesión actual). La fachada resuelve el handle en un solo lugar
(\code{SessionStore::resolve}) así toda operación hereda la
misma lógica de precedencia.

\begin{inthecode}
\filepath{graph\_hunter\_api/src/lib.rs} --- fachada
\code{GraphHunterApi}, \code{ApiState}.\\
\filepath{graph\_hunter\_api/src/state/session.rs} ---
\code{Session}, \code{SessionHandle}, \code{SessionStore}.
\end{inthecode}
